\section{Introduction}

The Basel problem asks for the sum of the reciprocal squares of the positive integers:
\begin{equation}
  S=\sum_{n=1}^{\infty}\frac{1}{n^{2}}
   =1+\frac{1}{2^{2}}+\frac{1}{3^{2}}+\frac{1}{4^{2}}+\cdots .
  \label{eq:basel}
\end{equation}
Mengoli posed it in 1644 and Euler evaluated the sum in 1735 as $\pi^{2}/6$. The left-hand
side is built out of integers alone, while the right-hand side is a circle constant, so one
expects a geometric reading of the identity.

Such a reading was given by W\"astlund \cite{wastlund} in a short note and popularised by
3Blue1Brown \cite{threeblueonebrown}. It uses nothing beyond Euclidean geometry: circles,
chords, right angles, and the inverse-square law for light. A lamp on a circle is a point
source, the common tangent point $P$ of two circles is the observer, and a term $1/d^{2}$ is
read as the brightness that a lamp at distance $d$ delivers at $P$. Doubling the circle
replaces every lamp by two equivalent ones on the larger circle, so the total brightness at
$P$ is preserved; in the limit the circle flattens into a line and its lamps fall on the odd
integers.

Normalise the construction so that the circle of circumference $2$ carries a single lamp, at
the point antipodal to $P$, at distance $2/\pi$. What the figure then delivers is the odd
terms.

\begin{theorem}[W\"astlund]\label{thm:wastlund}
With the normalisation above, the construction gives
\begin{equation}
  \sum_{k=0}^{\infty}\frac{1}{(2k+1)^{2}}=\frac{\pi^{2}}{8}.
  \label{eq:odd-sum}
\end{equation}
\end{theorem}

Getting $S$ out of Theorem \ref{thm:wastlund} costs one algebraic remark: the even terms are
$\frac14 S$, so $\frac34 S=\frac{\pi^{2}}{8}$. That remark is the only non-geometric step in
the whole argument, and it is the step this note removes.

\begin{theorem}[this note]\label{thm:direct}
Modify the doubling as follows: at every round, after the lamps have been replaced by their
images on the larger circle, add one lamp at the point of that circle antipodal to $P$. Then
every nonzero integer point of the limiting line carries a lamp, the brightness they deliver
at $P$ equals the brightness of the added lamps, and
\begin{equation}
  S=\frac{1}{2}\left(\frac{\pi^{2}}{2^{2}}+\frac{\pi^{2}}{4^{2}}+\frac{\pi^{2}}{8^{2}}+\cdots\right)
   =\frac{\pi^{2}}{6}.
  \label{eq:direct}
\end{equation}
The factor $\frac12$ is there because the line carries each distance twice, at $n$ and at $-n$:
what the picture delivers is $2S$. The term of round $k$ is the brightness $\pi^{2}/(2^{k+1})^{2}$
of the single lamp added in that round, which stands at distance $2^{k+1}/\pi$ from $P$.
\end{theorem}

Only one lamp has to be added each round, and it always sits in the same place: the antipodal
position is exactly the one the replacement cannot produce, since its preimage would be a lamp
at $P$ itself. Every other integer position is filled exactly once.

Equalities \eqref{eq:basel} and $\pi^{2}/6$ have been reached by many different routes;
Kalman's survey collects six of them \cite{kalman} and Matsuoka gives an elementary one
\cite{matsuoka}. Closest in spirit to the picture below is Dumont's heuristic argument,
made rigorous by Komornik, which ties the Gregory--Leibniz and Basel sums together
\cite{komornik}.

Every account of this geometric route that we were able to find has the same shape. W\"astlund's
own note \cite{wastlund}, the geometric-proof section of the Wikipedia article \cite{wikipedia},
and the video version by 3Blue1Brown \cite{threeblueonebrown} all carry the doubling to the
limit, read off the brightness of the lamps at the odd integers, and then recover $S$ from the
observation that the even terms are a quarter of it. As far as the author could find, no earlier
figure of this kind gives the whole series itself: each of these accounts stops at a sub-series
and closes with the same two lines of algebra, which is what \eqref{eq:direct} avoids. None of
them is the route we modify: what follows concerns only W\"astlund's construction
\cite{wastlund}.

Section \ref{sec:wastlund} reproduces W\"astlund's proof; Section \ref{sec:one-more-step}
contains the modification and the proof of Theorem \ref{thm:direct}, which is the only new
content of this note.
